% Options for packages loaded elsewhere
\PassOptionsToPackage{unicode}{hyperref}
\PassOptionsToPackage{hyphens}{url}
\documentclass[
]{article}
\usepackage{xcolor}
\usepackage[margin=1in]{geometry}
\usepackage{amsmath,amssymb}
\setcounter{secnumdepth}{-\maxdimen} % remove section numbering
\usepackage{iftex}
\ifPDFTeX
  \usepackage[T1]{fontenc}
  \usepackage[utf8]{inputenc}
  \usepackage{textcomp} % provide euro and other symbols
\else % if luatex or xetex
  \usepackage{unicode-math} % this also loads fontspec
  \defaultfontfeatures{Scale=MatchLowercase}
  \defaultfontfeatures[\rmfamily]{Ligatures=TeX,Scale=1}
\fi
\usepackage{lmodern}
\ifPDFTeX\else
  % xetex/luatex font selection
\fi
% Use upquote if available, for straight quotes in verbatim environments
\IfFileExists{upquote.sty}{\usepackage{upquote}}{}
\IfFileExists{microtype.sty}{% use microtype if available
  \usepackage[]{microtype}
  \UseMicrotypeSet[protrusion]{basicmath} % disable protrusion for tt fonts
}{}
\makeatletter
\@ifundefined{KOMAClassName}{% if non-KOMA class
  \IfFileExists{parskip.sty}{%
    \usepackage{parskip}
  }{% else
    \setlength{\parindent}{0pt}
    \setlength{\parskip}{6pt plus 2pt minus 1pt}}
}{% if KOMA class
  \KOMAoptions{parskip=half}}
\makeatother
\usepackage{color}
\usepackage{fancyvrb}
\newcommand{\VerbBar}{|}
\newcommand{\VERB}{\Verb[commandchars=\\\{\}]}
\DefineVerbatimEnvironment{Highlighting}{Verbatim}{commandchars=\\\{\}}
% Add ',fontsize=\small' for more characters per line
\usepackage{framed}
\definecolor{shadecolor}{RGB}{248,248,248}
\newenvironment{Shaded}{\begin{snugshade}}{\end{snugshade}}
\newcommand{\AlertTok}[1]{\textcolor[rgb]{0.94,0.16,0.16}{#1}}
\newcommand{\AnnotationTok}[1]{\textcolor[rgb]{0.56,0.35,0.01}{\textbf{\textit{#1}}}}
\newcommand{\AttributeTok}[1]{\textcolor[rgb]{0.13,0.29,0.53}{#1}}
\newcommand{\BaseNTok}[1]{\textcolor[rgb]{0.00,0.00,0.81}{#1}}
\newcommand{\BuiltInTok}[1]{#1}
\newcommand{\CharTok}[1]{\textcolor[rgb]{0.31,0.60,0.02}{#1}}
\newcommand{\CommentTok}[1]{\textcolor[rgb]{0.56,0.35,0.01}{\textit{#1}}}
\newcommand{\CommentVarTok}[1]{\textcolor[rgb]{0.56,0.35,0.01}{\textbf{\textit{#1}}}}
\newcommand{\ConstantTok}[1]{\textcolor[rgb]{0.56,0.35,0.01}{#1}}
\newcommand{\ControlFlowTok}[1]{\textcolor[rgb]{0.13,0.29,0.53}{\textbf{#1}}}
\newcommand{\DataTypeTok}[1]{\textcolor[rgb]{0.13,0.29,0.53}{#1}}
\newcommand{\DecValTok}[1]{\textcolor[rgb]{0.00,0.00,0.81}{#1}}
\newcommand{\DocumentationTok}[1]{\textcolor[rgb]{0.56,0.35,0.01}{\textbf{\textit{#1}}}}
\newcommand{\ErrorTok}[1]{\textcolor[rgb]{0.64,0.00,0.00}{\textbf{#1}}}
\newcommand{\ExtensionTok}[1]{#1}
\newcommand{\FloatTok}[1]{\textcolor[rgb]{0.00,0.00,0.81}{#1}}
\newcommand{\FunctionTok}[1]{\textcolor[rgb]{0.13,0.29,0.53}{\textbf{#1}}}
\newcommand{\ImportTok}[1]{#1}
\newcommand{\InformationTok}[1]{\textcolor[rgb]{0.56,0.35,0.01}{\textbf{\textit{#1}}}}
\newcommand{\KeywordTok}[1]{\textcolor[rgb]{0.13,0.29,0.53}{\textbf{#1}}}
\newcommand{\NormalTok}[1]{#1}
\newcommand{\OperatorTok}[1]{\textcolor[rgb]{0.81,0.36,0.00}{\textbf{#1}}}
\newcommand{\OtherTok}[1]{\textcolor[rgb]{0.56,0.35,0.01}{#1}}
\newcommand{\PreprocessorTok}[1]{\textcolor[rgb]{0.56,0.35,0.01}{\textit{#1}}}
\newcommand{\RegionMarkerTok}[1]{#1}
\newcommand{\SpecialCharTok}[1]{\textcolor[rgb]{0.81,0.36,0.00}{\textbf{#1}}}
\newcommand{\SpecialStringTok}[1]{\textcolor[rgb]{0.31,0.60,0.02}{#1}}
\newcommand{\StringTok}[1]{\textcolor[rgb]{0.31,0.60,0.02}{#1}}
\newcommand{\VariableTok}[1]{\textcolor[rgb]{0.00,0.00,0.00}{#1}}
\newcommand{\VerbatimStringTok}[1]{\textcolor[rgb]{0.31,0.60,0.02}{#1}}
\newcommand{\WarningTok}[1]{\textcolor[rgb]{0.56,0.35,0.01}{\textbf{\textit{#1}}}}
\usepackage{graphicx}
\makeatletter
\newsavebox\pandoc@box
\newcommand*\pandocbounded[1]{% scales image to fit in text height/width
  \sbox\pandoc@box{#1}%
  \Gscale@div\@tempa{\textheight}{\dimexpr\ht\pandoc@box+\dp\pandoc@box\relax}%
  \Gscale@div\@tempb{\linewidth}{\wd\pandoc@box}%
  \ifdim\@tempb\p@<\@tempa\p@\let\@tempa\@tempb\fi% select the smaller of both
  \ifdim\@tempa\p@<\p@\scalebox{\@tempa}{\usebox\pandoc@box}%
  \else\usebox{\pandoc@box}%
  \fi%
}
% Set default figure placement to htbp
\def\fps@figure{htbp}
\makeatother
\setlength{\emergencystretch}{3em} % prevent overfull lines
\providecommand{\tightlist}{%
  \setlength{\itemsep}{0pt}\setlength{\parskip}{0pt}}
\usepackage{bookmark}
\IfFileExists{xurl.sty}{\usepackage{xurl}}{} % add URL line breaks if available
\urlstyle{same}
\hypersetup{
  pdftitle={Credible Intervals},
  hidelinks,
  pdfcreator={LaTeX via pandoc}}

\title{Credible Intervals}
\author{}
\date{\vspace{-2.5em}2026-04-17}

\begin{document}
\maketitle

\subsection{Introduction}\label{introduction}

A credible interval is the Bayesian summary of posterior uncertainty
most often reported alongside a point estimate. Unlike a confidence
interval, whose probability statement refers to long-run frequency under
hypothetical resampling, a credible interval supports the direct claim
that ``given the data and prior, there is a 95 \% probability that the
parameter lies inside this interval.'' That single sentence is the
reason students consistently misinterpret confidence intervals --- and
the reason credible intervals are increasingly preferred in applied
biomedical reporting.

\subsection{Prerequisites}\label{prerequisites}

A working understanding of Bayesian inference, posterior distributions,
and the distinction between a posterior and a sampling distribution.

\subsection{Theory}\label{theory}

Two definitions dominate practice:

\begin{itemize}
\tightlist
\item
  \textbf{Equal-tailed interval (ETI)}: the interval bounded by the
  \(\alpha/2\) and \(1 - \alpha/2\) quantiles of the posterior. It is
  invariant under monotone reparameterisation and easy to compute from
  MCMC draws.
\item
  \textbf{Highest posterior density (HPD) interval}: the narrowest set
  with \(1 - \alpha\) posterior mass; for unimodal posteriors it is a
  single interval centred on the region of highest density, and it is
  always at least as short as the ETI.
\end{itemize}

For symmetric posteriors the two intervals coincide; for skewed
posteriors they can differ markedly. A 95 \% level is conventional but
arbitrary --- some Bayesian methodologists advocate 89 \% to discourage
misreading the threshold as a hypothesis test.

\subsection{Assumptions}\label{assumptions}

A proper posterior distribution and either MCMC draws (post-warm-up,
with adequate effective sample size) or a closed-form posterior. ETI
endpoints are quantiles and tolerate small samples better than HPD
endpoints, which depend on density estimation in the tails.

\subsection{R Implementation}\label{r-implementation}

\begin{Shaded}
\begin{Highlighting}[]
\FunctionTok{library}\NormalTok{(HDInterval)}

\FunctionTok{set.seed}\NormalTok{(}\DecValTok{2026}\NormalTok{)}
\NormalTok{theta }\OtherTok{\textless{}{-}} \FunctionTok{rbeta}\NormalTok{(}\DecValTok{10000}\NormalTok{, }\DecValTok{10}\NormalTok{, }\DecValTok{4}\NormalTok{)}

\CommentTok{\# Equal{-}tailed}
\FunctionTok{quantile}\NormalTok{(theta, }\FunctionTok{c}\NormalTok{(}\FloatTok{0.025}\NormalTok{, }\FloatTok{0.975}\NormalTok{))}

\CommentTok{\# HPD}
\FunctionTok{hdi}\NormalTok{(theta, }\AttributeTok{credMass =} \FloatTok{0.95}\NormalTok{)}
\end{Highlighting}
\end{Shaded}

\subsection{Output \& Results}\label{output-results}

\texttt{quantile()} returns the ETI; \texttt{hdi()} returns the HPD. For
a Beta(10, 4) posterior the ETI is approximately \((0.50, 0.90)\) and
the HPD is approximately \((0.53, 0.93)\) --- slightly shifted toward
the right because the posterior is left-skewed and the HPD chases the
highest density region rather than splitting tail probability evenly.

\subsection{Interpretation}\label{interpretation}

A reporting sentence: ``The posterior probability had a 95 \%
equal-tailed credible interval of 0.50 to 0.90; the
highest-posterior-density interval was 0.53 to 0.93. There is a 95 \%
posterior probability that the true proportion falls in either interval,
given the prior and the observed data.'' Always state which interval is
reported; mixing definitions across tables in the same paper is a common
reproducibility headache.

\subsection{Practical Tips}\label{practical-tips}

\begin{itemize}
\tightlist
\item
  Report ETIs for symmetric or near-symmetric posteriors; they are more
  familiar and reparameterisation-invariant.
\item
  Report HPDs for skewed posteriors when the additional precision
  matters and the density estimate is stable.
\item
  HPD intervals can be disjoint under multimodal posteriors; set
  \texttt{allowSplit\ =\ TRUE} in \texttt{HDInterval::hdi()} so the
  structure is preserved.
\item
  Disclose the prior in any reporting sentence; the credible interval is
  not a property of the data alone.
\item
  Use post-warm-up MCMC draws and check ESS in the tails ---
  credible-interval endpoints depend on quantile estimation, which needs
  more effective draws than the mean.
\end{itemize}

\subsection{R Packages Used}\label{r-packages-used}

\texttt{HDInterval} for the canonical \texttt{hdi()} function,
\texttt{bayestestR::ci()} for an integrated wrapper that selects between
ETI and HPD on \texttt{brms} and \texttt{rstanarm} model objects, and
base R (\texttt{quantile()}) for the equal-tailed case.

\subsection{For Reviewers}\label{for-reviewers}

\textbf{What to look for in a paper using this method.}

\begin{itemize}
\tightlist
\item
  \textbf{Common misapplications.}
\item
  \textbf{Diagnostics that should be reported but often aren't.}
\item
  \textbf{Red flags in tables and figures.}
\item
  \textbf{What to verify.}
\item
  \textbf{What an adequate Methods paragraph must contain.}
\end{itemize}

\subsection{See also --- labs in this
chapter}\label{see-also-labs-in-this-chapter}

\begin{itemize}
\tightlist
\item
  \href{labs/lab-bayesian-thinking-control-vs-decision-error.Rmd}{Bayesian
  thinking; α control vs decision error}
\item
  \href{labs/lab-brms-stan-loo-hierarchical-models.Rmd}{brms / Stan,
  LOO, hierarchical models}
\end{itemize}

\end{document}
